
\documentclass[11pt,a4paper]{article}

\usepackage[utf8]{inputenc}
\usepackage[T1]{fontenc}
\usepackage[french]{babel}
\usepackage{amsmath,amssymb,amsfonts}
\usepackage{geometry}
\usepackage{graphicx}
\usepackage{booktabs}
\usepackage{hyperref}
\usepackage{array}
\usepackage{physics}

\geometry{margin=2.5cm}

\title{
\textbf{BuP Paper 10 -- Limite quasi-newtonienne du propagateur d'intrication}\\
\large Fonction de Green discrète, dimension spectrale et corrections BuP
}

\author{Farid Hamdad}
\date{2026}

\begin{document}

\maketitle

\begin{abstract}
Dans le programme Bottom-Up Quantum Gravity (BuP), la géométrie et la gravité
émergent de la structure d'intrication d'un état quantique fini. Les articles
précédents ont introduit un graphe d'information mutuelle, un Laplacien
d'intrication \(L_{\rm ent}\), une dimension spectrale effective \(d_s\), ainsi
qu'une source tensorielle candidate de matière émergente.

Le présent article étudie la pseudo-inverse du Laplacien,
\[
L_{\rm ent}^{+},
\]
comme propagateur gravitationnel discret. Nous testons la réponse à une source
ponctuelle,
\[
L_{\rm ent}\Phi
=
\delta_{i_0}
-
\frac{1}{N},
\]
et analysons le profil radial du potentiel effectif,
\[
\Phi(r)
\simeq
C
+
A r^{-\alpha}.
\]

Nous montrons d'abord, sur des graphes géométriques contrôlés, que le protocole
point-source permet d'extraire une loi de Green en dimension trois. Nous
construisons ensuite des matrices d'information mutuelle géométriques avec
\(N=800\), \(\lambda=0.3\), \(k=12\), dont la dimension spectrale mesurée par
noyau de chaleur est proche de trois.

Pour la réalisation principale, nous obtenons
\[
d_s=2.964,
\qquad
\alpha_{\rm pred}=d_s-2=0.964.
\]
Dans la fenêtre intermédiaire
\[
r/r_{\max}\in[0.15,0.50],
\]
le fit inter-centres du propagateur donne
\[
\alpha_{\rm med}=1.1046
\]
pour le filtre strict \(R^2>0.85\). Des vérifications avec davantage de centres
et un autre seed de sélection donnent respectivement
\[
\alpha_{\rm med}=1.1547
\quad \text{et} \quad
\alpha_{\rm med}=1.2022.
\]
Une seconde réalisation indépendante de la matrice MI donne
\[
d_s=2.970,
\qquad
\alpha_{\rm med}=1.2597.
\]

Ces résultats établissent l'existence d'une fenêtre quasi-newtonienne du
propagateur d'intrication. La loi de Newton \(1/r\) n'est pas imposée dans BuP :
elle apparaît comme une limite effective lorsque le graphe atteint un régime
spectral proche de trois dimensions, avec des corrections discrètes et
spectrales mesurables.
\end{abstract}

\tableofcontents

\section{Introduction}

Le programme BuP repose sur l'hypothèse que la géométrie n'est pas fondamentale.
Elle émerge d'une structure d'intrication encodée par une matrice d'information
mutuelle,
\[
W_{ij}=I(i:j).
\]
À partir de cette matrice, on reconstruit un graphe pondéré, puis un Laplacien
d'intrication,
\[
L_{\rm ent}.
\]

Les travaux précédents ont montré que ce Laplacien permet de définir une
dimension spectrale effective,
\[
d_s(\tau)
=
-2
\frac{d\log Z(\tau)}{d\log \tau},
\qquad
Z(\tau)=\Tr e^{-\tau L_{\rm ent}},
\]
et que les excitations locales du graphe peuvent être interprétées comme des
sources candidates de matière.

Le problème étudié ici est le suivant : si \(L_{\rm ent}\) joue le rôle d'un
opérateur géométrique discret, alors sa pseudo-inverse
\[
L_{\rm ent}^{+}
\]
doit jouer le rôle d'un propagateur de Green. Dans une géométrie continue
tridimensionnelle, la fonction de Green du Laplacien donne le potentiel
newtonien,
\[
\Phi(r)\sim \frac{1}{r}.
\]
La question centrale devient donc :
\[
\boxed{
\text{la loi newtonienne peut-elle émerger du propagateur }L_{\rm ent}^{+}?
}
\]

\section{Propagateur d'intrication}

\subsection{Définition}

Soit un graphe d'intrication pondéré par une matrice \(W_{ij}\). Le Laplacien
combinatoire associé est
\[
L_{\rm ent}
=
D-W,
\]
où
\[
D_{ii}
=
\sum_j W_{ij}.
\]

La pseudo-inverse de Moore--Penrose de ce Laplacien est notée
\[
L_{\rm ent}^{+}.
\]

Comme le Laplacien possède un mode nul constant,
\[
L_{\rm ent}\mathbf{1}=0,
\]
la réponse à une source ponctuelle doit être rendue compatible avec la
condition de moyenne nulle. Nous résolvons donc :
\[
L_{\rm ent}\Phi
=
\delta_{i_0}
-
\frac{1}{N}\mathbf{1}.
\]

La solution est
\[
\Phi
=
L_{\rm ent}^{+}
\left(
\delta_{i_0}
-
\frac{1}{N}\mathbf{1}
\right).
\]

\subsection{Profil radial}

Pour un centre \(i_0\), on définit une distance d'intrication \(r(i_0,j)\), puis
on ajuste le potentiel par :
\[
\Phi(r)
=
C
+
A r^{-\alpha}.
\]

L'exposant \(\alpha\) caractérise la loi effective du propagateur.

Dans une limite tridimensionnelle continue, on attend :
\[
\alpha\simeq1.
\]

Plus généralement, une loi spectrale naturelle est :
\[
\alpha
\simeq
d_s-2.
\]

Cette relation est exacte pour le Laplacien continu dans des géométries
régulières. Sur un graphe fini et irrégulier, elle doit être comprise comme une
relation asymptotique à tester numériquement.

\section{Calibration sur graphes géométriques contrôlés}

\subsection{Motivation}

Avant de tester des matrices d'information mutuelle BuP, il faut vérifier que
le protocole point-source est capable de retrouver une loi de Green correcte
sur des graphes géométriques connus.

Nous considérons des grilles 1D, 2D et 3D. Pour chaque graphe, nous résolvons :
\[
L\Phi
=
\delta_{i_0}
-
\frac{1}{N}.
\]

Le profil est ensuite ajusté par :
\[
\Phi(r)=C+A r^{-\alpha}.
\]

\subsection{Résultats}

Les tests sur grilles contrôlées montrent que la réponse ponctuelle produit une
loi de puissance propre en dimension trois.

Par exemple, pour la grille 3D \(L=11\), on obtient :
\[
\alpha=1.152,
\qquad
R^2=0.996.
\]

Pour \(L=13\), le fit reste extrêmement propre :
\[
\alpha=1.330,
\qquad
R^2=0.9998.
\]

Ces résultats montrent que le protocole point-source extrait bien une loi de
Green discrète. La valeur exacte de \(\alpha\) dépend de la taille finie, des
bords, de la fenêtre de fit et de la discrétisation.

\subsection{Conclusion de la calibration}

Le protocole retenu n'est pas la moyenne globale de toutes les entrées de
\(L^+\). Cette moyenne mélange les signes, les effets de bord et les paires non
équivalentes.

Le bon protocole physique est la réponse à une source ponctuelle :
\[
L\Phi=\delta_{i_0}-\frac{1}{N}.
\]

Cette réponse fournit une fonction de Green discrète exploitable.

\section{Comparaison du modèle newtonien et du modèle libre}

\subsection{Deux modèles}

Sur les graphes 3D contrôlés, nous comparons deux modèles.

Le modèle newtonien fixé :
\[
\Phi_{\rm Newt}(r)
=
C+\frac{A}{r}.
\]

Le modèle à exposant libre :
\[
\Phi_{\rm free}(r)
=
C+A r^{-\alpha}.
\]

La comparaison utilise \(R^2\), AIC et BIC.

\subsection{Interprétation}

Le modèle newtonien donne généralement un excellent ajustement, mais le modèle
libre est parfois préféré. Cela indique que le propagateur discret possède une
loi quasi-newtonienne, avec des corrections de réseau mesurables.

La conclusion n'est donc pas :
\[
\alpha=1
\quad
\text{exactement}.
\]

La conclusion correcte est :
\[
\boxed{
\alpha\simeq1
\quad
\text{dans une fenêtre intermédiaire.}
}
\]

\section{Construction d'une matrice MI quasi tridimensionnelle}

\subsection{Génération}

Nous générons des matrices MI géométriques de la forme :
\[
W_{ij}
=
\exp\left[
-
\left(
\frac{r_{ij}}{\lambda}
\right)^p
\right],
\]
où \(r_{ij}\) est une distance dans un espace ambiant 3D.

À partir de \(W_{ij}\), nous construisons un graphe \(k\)-NN et mesurons la
dimension spectrale par noyau de chaleur :
\[
Z(\tau)=\Tr e^{-\tau L_{\rm ent}},
\]
\[
d_s(\tau)
=
-2
\frac{d\log Z(\tau)}{d\log\tau}.
\]

\subsection{Candidat principal}

Le candidat principal utilise :
\[
N=800,
\qquad
\lambda=0.3,
\qquad
k=12.
\]

Pour la réalisation \(seed=0\), la dimension spectrale mesurée est :
\[
d_s=2.964.
\]

La relation spectrale attendue donne :
\[
\alpha_{\rm pred}
=
d_s-2
=
0.964.
\]

Ce graphe est donc très proche d'un régime spectral tridimensionnel.

\subsection{Deuxième réalisation}

Pour une seconde réalisation indépendante de la matrice MI, avec \(seed=1\),
nous obtenons :
\[
d_s=2.970,
\qquad
\alpha_{\rm pred}=0.970.
\]

Cela confirme que le régime quasi tridimensionnel ne dépend pas d'une unique
réalisation de la matrice.

\section{Test du propagateur sur matrice MI quasi-3D}

\subsection{Protocole}

Pour la matrice MI \(N=800,\lambda=0.3,k=12\), nous résolvons le problème
point-source :
\[
L_{\rm ent}\Phi_{i_0}
=
\delta_{i_0}
-
\frac{1}{N}.
\]

Le calcul est répété sur un sous-échantillon de centres \(i_0\). Pour chaque
centre, on ajuste :
\[
\Phi_{i_0}(r)
=
C
+
A r^{-\alpha_{i_0}}.
\]

Les centres sont ensuite filtrés par qualité de fit :
\[
R^2>0.70
\]
et
\[
R^2>0.85.
\]

\subsection{Importance de la fenêtre de fit}

Nous testons plusieurs fenêtres radiales normalisées :
\[
r/r_{\max}\in[0.10,0.40],
\]
\[
r/r_{\max}\in[0.15,0.50],
\]
\[
r/r_{\max}\in[0.20,0.60].
\]

Les résultats montrent que la fenêtre courte \([0.10,0.40]\) reste trop proche
du régime UV et donne un exposant sous-newtonien :
\[
\alpha_{\rm med}=0.5333
\quad
(R^2>0.85).
\]

La fenêtre centrale \([0.15,0.50]\) donne le meilleur accord avec la limite
quasi-newtonienne :
\[
\alpha_{\rm med}=1.1046
\quad
(R^2>0.85).
\]

La fenêtre plus large \([0.20,0.60]\) donne :
\[
\alpha_{\rm med}=1.1397
\quad
(R^2>0.85).
\]

\subsection{Résultat principal}

Dans la fenêtre centrale :
\[
r/r_{\max}\in[0.15,0.50],
\]
avec \(150\) centres et le filtre strict \(R^2>0.85\), on obtient :
\[
n=109,
\]
\[
\alpha_{\rm med}=1.1046,
\]
\[
\alpha_{\rm mean}=1.1500,
\]
\[
\sigma_\alpha=0.3285,
\]
\[
\Delta
=
\alpha_{\rm med}
-
(d_s-2)
=
0.1406.
\]

Ce résultat est proche de :
\[
\alpha_{\rm pred}=0.964.
\]

Il est aussi proche de la loi newtonienne :
\[
\alpha=1.
\]

\section{Tests de robustesse}

\subsection{Davantage de centres}

Nous répétons le test principal avec \(300\) centres, même fenêtre
\([0.15,0.50]\), même matrice MI.

Pour \(R^2>0.85\), nous obtenons :
\[
n=202,
\]
\[
\alpha_{\rm med}=1.1547,
\]
\[
\Delta=0.1906.
\]

\subsection{Autre seed de sélection des centres}

Nous répétons le test principal avec \(150\) centres mais un autre seed de
sélection.

Pour \(R^2>0.85\), nous obtenons :
\[
n=102,
\]
\[
\alpha_{\rm med}=1.2022,
\]
\[
\Delta=0.2381.
\]

\subsection{Autre seed de matrice MI}

Nous générons une deuxième réalisation indépendante de la matrice MI avec les
mêmes paramètres :
\[
N=800,
\qquad
\lambda=0.3,
\qquad
k=12,
\]
mais avec \(seed=1\).

Cette réalisation donne :
\[
d_s=2.970,
\qquad
d_s-2=0.970.
\]

Dans la même fenêtre \([0.15,0.50]\), pour \(R^2>0.85\), nous obtenons :
\[
n=104,
\]
\[
\alpha_{\rm med}=1.2597,
\]
\[
\Delta=0.2893.
\]

\subsection{Conclusion des tests}

Les tests donnent :
\[
\alpha_{\rm med}
\simeq
1.10-1.26.
\]

La valeur exacte varie selon l'échantillon de centres et la réalisation de la
matrice MI, mais le régime reste quasi-newtonien.

\begin{table}[h]
\centering
\begin{tabular}{lccccc}
\toprule
Test & Fenêtre & Centres & Source & Filtre & \(\alpha_{\rm med}\) \\
\midrule
Principal & [0.15,0.50] & 150 & matrix seed0 & \(R^2>0.85\) & 1.1046 \\
Plus de centres & [0.15,0.50] & 300 & matrix seed0 & \(R^2>0.85\) & 1.1547 \\
Autre seed centres & [0.15,0.50] & 150 & matrix seed0 & \(R^2>0.85\) & 1.2022 \\
Autre matrice MI & [0.15,0.50] & 150 & matrix seed1 & \(R^2>0.85\) & 1.2597 \\
\bottomrule
\end{tabular}
\caption{Robustesse de l'exposant quasi-newtonien sur les matrices MI \(N=800\).}
\end{table}

\section{Interprétation physique}

\subsection{Newton comme limite effective}

La loi de Newton :
\[
\Phi(r)
\sim
\frac{1}{r}
\]
est une loi idéale du Laplacien continu en dimension trois.

Dans BuP, on ne part pas d'un espace continu. On part d'un graphe fini issu
d'une structure d'intrication. La loi newtonienne ne doit donc pas être imposée
à la main. Elle doit émerger dans une fenêtre d'échelle.

Les résultats de ce papier montrent que, lorsque :
\[
d_s\simeq3,
\]
le propagateur d'intrication possède effectivement une fenêtre où :
\[
\alpha\simeq1.
\]

\subsection{Correction BuP}

L'exposant mesuré n'est pas exactement :
\[
\alpha=1.
\]

Il est plutôt :
\[
\alpha_{\rm med}\simeq1.1-1.3.
\]

Cette déviation peut être interprétée comme une correction de graphe fini,
liée à :

\begin{itemize}
\item la dimension spectrale effective \(d_s\),
\item la taille finie \(N\),
\item l'hétérogénéité locale,
\item la fenêtre de fit,
\item la discrétisation du propagateur.
\end{itemize}

Ainsi, Paper 10 ne prétend pas démontrer Newton exact. Il établit :
\[
\boxed{
\text{une fenêtre quasi-newtonienne robuste du propagateur }L_{\rm ent}^{+}.
}
\]

\section{Lien avec la dimension critique cosmologique}

Dans les travaux cosmologiques BuP, une dimension critique proche de
\[
D_C\simeq3.06
\]
a été identifiée.

Si la relation asymptotique
\[
\alpha\simeq d_s-2
\]
s'applique dans la limite continue, alors :
\[
\alpha_{\rm BuP}
\simeq
D_C-2
\simeq
1.06.
\]

Cette valeur est proche des exposants mesurés dans la fenêtre quasi-newtonienne :
\[
\alpha_{\rm med}\simeq1.10-1.26.
\]

Nous n'interprétons pas encore cette correspondance comme une preuve. Nous la
présentons comme une conjecture testable :
\[
\boxed{
\alpha_\infty
\simeq
D_C-2.
}
\]

Cette conjecture relierait directement :

\[
\text{dimension spectrale cosmologique}
\longrightarrow
\text{loi gravitationnelle effective}.
\]

\section{Limites}

Les résultats restent numériques et doivent être consolidés.

Les principales limites sont :

\begin{enumerate}
\item La relation \(\alpha=d_s-2\) n'est pas encore démontrée analytiquement sur
graphe discret général.
\item La valeur de \(\alpha\) dépend de la fenêtre de fit.
\item La matrice MI utilisée est une matrice géométrique contrôlée, non encore
une matrice MI quantique complète issue d'un état \(\ket{\Psi}\).
\item Les effets de taille finie restent importants.
\item Le passage aux observations astrophysiques nécessite un modèle dynamique
complet.
\end{enumerate}

\section{Perspectives}

Paper 11 devra viser trois objectifs.

Premièrement, fournir une démonstration analytique reliant :
\[
K(\tau,x,x)
\sim
\tau^{-d_s/2}
\]
à :
\[
L^{+}(r)
\sim
r^{-(d_s-2)}.
\]

Deuxièmement, tester la convergence grand \(N\) :
\[
\alpha(N)
\to
d_s(N)-2.
\]

Troisièmement, relier la correction :
\[
\alpha-1
\]
aux observables galactiques et cosmologiques.

\section{Conclusion}

Nous avons montré que la pseudo-inverse du Laplacien d'intrication,
\[
L_{\rm ent}^{+},
\]
se comporte comme un propagateur de Green discret.

Sur des graphes contrôlés, le protocole point-source extrait une loi de
puissance en dimension trois. Sur des matrices MI géométriques quasi
tridimensionnelles avec :
\[
N=800,
\qquad
k=12,
\qquad
\lambda=0.3,
\qquad
d_s\simeq2.96-2.97,
\]
nous obtenons dans la fenêtre intermédiaire :
\[
r/r_{\max}\in[0.15,0.50],
\]
des exposants médians :
\[
\alpha_{\rm med}\simeq1.10-1.26
\quad
(R^2>0.85).
\]

Ainsi, la loi newtonienne n'est pas postulée dans BuP. Elle apparaît comme une
approximation effective du propagateur d'intrication lorsque le graphe atteint
un régime spectral proche de trois dimensions.

Le résultat central de Paper 10 est donc :
\[
\boxed{
d_s\simeq3
\quad
\Longrightarrow
\quad
L_{\rm ent}^{+}
\text{ possède une fenêtre quasi-newtonienne.}
}
\]

\end{document}
